\section{Background and Related Work}

Electronic limit order markets are usually organised as continuous double
auctions: the exchange maintains a bid book and an ask book, and matches orders
when the best bid is at least the best ask. In the common price-time-priority
model, higher bids outrank lower bids, lower asks outrank higher asks, and
orders at the same price are filled in arrival order~\cite{harris}. This basic
rule is widely used in practice, although real venues often add richer order
types, recovery protocols, and venue-specific priority overlays.

That gap between a clean matching rule and a production exchange motivates the
scope of this paper. Real matching engines must also satisfy low-latency,
fault-tolerance, auditability, and regulatory requirements, but those concerns
are not the main target here. Instead, this work isolates the exchange core:
a single-asset limit order book with explicit state transitions for order
submission and matching. The goal is not to model full market microstructure,
but to capture the central correctness question in a form suitable for formal
verification.

Methodologically, the project is closest to work on transition-system
verification and verified stateful systems. Veil is a natural fit because it
supports transition-system modelling, model checking on small instances, and
automated invariant checking for safety properties~\cite{veilpost}. Lean is
useful as a supporting tool for lemmas about sorted data structures and
arithmetic, such as preservation of ordered books and conservation of
quantities. If extended to an executable reference implementation, Velvet
offers a way to connect an algorithmic matching loop back to the abstract
specification~\cite{velvetpost}. Loom is also relevant as part of the broader
Verse-Lab ecosystem around foundational multi-modal verification~\cite{loom}.
